\begin{claim}{F.28}{Proved}{All finite affine marginal images on the three-parameter product-one locus}
Retain the actual bases $Y_i^-=X_i$ and
$Y_i^+=\mathcal B_+(X_i)$ of F.27, with
$\mathcal B_+$ the based comparison of F.23.13.
For every finite field $k$ and continuous
$\rho:V_\pm\to\operatorname{Aff}_1(k)$ with scalar
character $\chi$, the locus
\[
 \chi(Y_4^\pm)=\chi(Y_5^\pm)=\chi(Y_6^\pm)=1,
 \qquad
 \chi(Y_1^\pm)\chi(Y_2^\pm)\chi(Y_3^\pm)
                         \chi(Y_7^\pm)=1
                                                               \tag{F.28.1}
\]
has the exact image equality
\[
                         \rho(\mathcal K_\pm)=\rho(R_\pm).
                                                               \tag{F.28.2}
\]
This includes every characteristic, every finite
scalar-field extension degree and every larger
original coefficient field.

For a negative character
$(r,s,t,1,1,1,c^{-1})$, $c=rst\ne1$, put
$k_1=\mathbb F_p(c)$ and $k_0=\mathbb F_p(r,s,t)$.
In the case $k_1\subsetneq k_0$, let
$q_1=|k_1|$, $m=[k_0:k_1]$ and
$\chi_j=\chi^{q_1^j}$, $0\leq j<m$.
The summed continuous $H^1$ restriction for these
$m$ characters is injective, and its cocycle
kernel is exactly the principal tuples with
parameter sum zero. Its rank is therefore
$6m+1$ on $7m$ cocycle coordinates.
\end{claim}

\begin{proof}
\emph{The actual fields and the already proved cases.}
Use the negative groups $K\leq R\leq V$ first.
The deleted generators $X_4,\ldots,X_7$ and the
actual normal-generation maps of F.21--F.22 give
\[
 \begin{gathered}
 C=\chi(K)=\chi(R)=\langle c^{-1}\rangle,
 \qquad D=\chi(V)=\langle r,s,t\rangle,\\
 k_1=\mathbb F_p(c),\qquad
 k_0=\mathbb F_p(r,s,t)=k_1(r,t),\qquad c=rst.
 \end{gathered}                                                \tag{F.28.3}
\]
The last equality follows from $s=c/(rt)$;
all character values are nonzero.
For $c=1$, the $R$ image is an abelian translation
group, and actual normal generation gives (F.28.2).
For $c\ne1$ and $k_0=k_1$, F.25 gives it.
For $r=1$, the character is exactly
$(1,s,t,1,1,1,(st)^{-1})$ in the unchanged actual
basis, so F.27 gives both the image equality
and the stated summed-restriction result.
At the coefficient-ring level this last comparison
is the homomorphism $r\mapsto1$, $s\mapsto s$,
$t\mapsto t$, sending $1-rst$ to $1-st$ in
the corresponding localized Laurent rings.
Thus it respects the precise maps used in F.27.

It remains to treat $c\ne1$, $r\ne1$ and
$k_1\subsetneq k_0$. These conditions imply
only the nonzero factors actually used below:
$r,s,t,c,c-1,r-1$. No condition on $s-1$,
$t-1$, $r+1$ or the characteristic is imposed.

\emph{The exact invariant inclusion and quotient maps.}
The literal coefficient homomorphism from F.24 is
\[
 \begin{aligned}
 \mathbb Z[t_1^{\pm1},\ldots,t_7^{\pm1}]&\longrightarrow
                   A=\mathbb Z[r^{\pm1},s^{\pm1},t^{\pm1}],\\
 (t_1,t_2,t_3,t_4,t_5,t_6,t_7)
          &\longmapsto(r,s,t,1,1,1,(rst)^{-1}).
 \end{aligned}                                                \tag{F.28.4}
\]
A monomial with exponent vector $(a_i)$ maps to
$r^{a_1-a_7}s^{a_2-a_7}t^{a_3-a_7}$.
These are identities on the specified discrete
free words. Their finite-field cocycle evaluations
extend uniquely to $V=\widehat F_7$ by the finite
affine homomorphisms, exactly as in F.27.7--8.

Let $\delta=(r,s,t,1,1,1,c^{-1})^{\mathsf T}
-\mathbf1$, with $\delta_7=(1-c)/c\ne0$.
The actual continuous cohomology space
$Q_\chi=H^1_{\rm cont}(V,k_{0,\chi})$ is $k_0^6$
through the projection and section
\[
 \begin{aligned}
 P(v)&=(v_i-\delta_i\delta_7^{-1}v_7)_{i=1}^6,
       &S(w)&=(w_1,\ldots,w_6,0)^{\mathsf T},\\
 PS&=I_6,
       &SP(v)&=v-\delta\delta_7^{-1}v_7.
 \end{aligned}                                                \tag{F.28.5}
\]
For the eleven actual F.23 source automorphisms
$B_l$, $0\leq l\leq10$, both directions preserve
$K$ and fix the character, as in F.27.
Let $A_l$ be their
specialized cocycle matrices. Then $A_l\delta=\delta$
and $\overline A_l=PA_lS$ are the actual source
operators on $Q_\chi$.

The boundary word $P_0=X_1\cdots X_7$ has
character one and derivative
$(1,r,rs,c,c,c,c)$. Pure Artin monodromy fixes
it, giving the quotient covector
$b=(1,r,rs,c,c,c)$.
The character-one meridians $X_4,X_5,X_6$ give
the three quotient covectors $e_4^*,e_5^*,e_6^*$,
as in F.27: the scalar of their actual conjugator
multiplies their cocycle value.
The four maps together define
\[
 \begin{aligned}
 \pi:Q_\chi&\longrightarrow k_0^4,
       &\pi(v)&=(bv,v_4,v_5,v_6)^{\mathsf T},\\
 \iota:k_0^2&\longrightarrow Q_\chi,
       &\iota(x,y)&=(-rx-rsy,x,y,0,0,0)^{\mathsf T}.
 \end{aligned}                                                \tag{F.28.6}
\]
Their kernel-image identity is
$\ker\pi=\operatorname{im}\iota=:W_\chi$.
The coordinate projection $(v_2,v_3)$ is inverse
to $\iota$ on this subspace. A coefficient section
of $\pi$ is
\[
 \sigma(z_\partial,z_4,z_5,z_6)
       =(z_\partial-c(z_4+z_5+z_6),0,0,z_4,z_5,z_6)^{\mathsf T}.
                                                               \tag{F.28.7}
\]
One has $\pi\sigma=I_4$ and
$\iota(v_2,v_3)+\sigma\pi(v)=v$.
Thus the basis matrix with columns
\[
 \iota(1,0),\quad\iota(0,1),\quad e_1,
             \quad e_4-ce_1,\quad e_5-ce_1,\quad e_6-ce_1
                                                               \tag{F.28.8}
\]
has inverse $v\mapsto(v_2,v_3,bv,v_4,v_5,v_6)$.
The displayed identities prove both inverse
coefficient maps explicitly.

In this basis the actual source matrices have form
\[
 \begin{pmatrix}G_l&E_l\\0&D_l\end{pmatrix},\qquad
 D_l=\operatorname{diag}
  (\lambda_{\partial,l},\lambda_{4,l},
                         \lambda_{5,l},\lambda_{6,l}).
                                                               \tag{F.28.9}
\]
Equivalently,
$\overline A_l\iota=\iota G_l$ and
$\pi\overline A_l=D_l\pi$.
All extension matrices $E_l$ are retained in
the certificate. These maps give the actual
invariant subspace $W_\chi$ and the four actual
one-dimensional quotient coordinates.

The literal source-word calculation gives the
following comparison table:
\[
 \begin{array}{c|c|cccc}
 l&G_l&\lambda_{\partial,l}&\lambda_{4,l}
                         &\lambda_{5,l}&\lambda_{6,l}\\\hline
 0&rI_2&1&r&r&r\\
 7&tI_2&t^{-1}&t&t&ct\\
 8&I_2&1&c&1&t^{-1}\\
 9&I_2&1&t^{-1}&c&1
 \end{array}
 \qquad \lambda_{\partial,2}=rt.
                                                               \tag{F.28.10}
\]
Every source index here is the original index
beginning at zero: $0=\texttt{old0}$,
$2=\texttt{old\_1\_1}$,
$7=\texttt{old\_-1\_1}$,
$8=\texttt{old\_-1\_2}$ and
$9=\texttt{old\_-1\_3}$.
The independent verifier checks all entries
against these exact labelled source words.

There are also explicit actual $K$ evaluations
on $W_\chi$. Let $a=u_{\texttt{cross\_4\_-1\_1}}$
be seed $25$, with $\chi(a)=c^{-1}$.
For $h_2=B_2(a)a^{-1}$ and $h_7=B_7(a)a^{-1}$,
both characters are one. A representative with
generator values $S\iota(x,y)$ satisfies
\[
 \bigl(d(h_2),d(h_7)\bigr)=((1-c)x,s(c-1)y).
                                                               \tag{F.28.11}
\]
Both evaluations therefore descend to cohomology.
Their inverse coefficient map is
$(u,v)\mapsto((1-c)^{-1}u,[s(c-1)]^{-1}v)$.
The determinant $-s(c-1)^2$ is nonzero in every
remaining case. These identities are also
verified from the actual seed and source words.

\emph{Frobenius transport and the classification needed for comparison.}
For $0\leq j<m$, put
$r_j=r^{q_1^j}$, $s_j=s^{q_1^j}$ and
$t_j=t^{q_1^j}$. Then $r_js_jt_j=c$.
The semilinear coefficient map
\[
             d\longmapsto d^{q_1^j},\qquad
             k_{0,\chi}\longrightarrow k_{0,\chi_j}
                                                               \tag{F.28.12}
\]
and its inverse, the power $q_1^{m-j}$ with the
identity at $j=0$, preserve cocycles, principals
and continuity. They commute with all actual
source precomposition maps. Consequently all
the preceding inclusions, projections, inverse
bases and comparison entries transport to
$Q_j=H^1_{\rm cont}(V,k_{0,\chi_j})$ and its
invariant subspace $W_j$ by the substitutions
$(r,s,t)\mapsto(r_j,s_j,t_j)$.
The characters agree on $K$, since
$\chi(K)=C\subset k_1^\times$.

The pairs $(r_j,t_j)$ are distinct. Equality of
two pairs would make a nonidentity $k_1$-Frobenius
power fix the generators of $k_0=k_1(r,t)$.
It would then fix every element of $k_0$;
the nonzero polynomial $X^{q_1^a}-X$, with
$0<a<m$, has degree smaller than $|k_0|$ and
cannot vanish on all those elements. Thus
\[
                  (r_j,t_j)=(r_l,t_l)\quad\Longrightarrow\quad j=l.
                                                               \tag{F.28.13}
\]
Also every $r_j\ne1$, since Frobenius is a field
automorphism and $r\ne1$.

Consider a nonzero invariant subspace $U\leq Q_j$
having no proper nonzero invariant subspace.
Its intersection with $W_j$ is invariant.
If the intersection is nonzero, $U\leq W_j$.
If it is zero, $\pi_j|U$ is injective. At least
one of its four actual coordinate projections
is nonzero. That projection has invariant
kernel, hence is injective by the defining
property of $U$. Thus $U$ is one-dimensional
and its source character is one of
$\lambda_{\partial,j},\lambda_{4,j},
\lambda_{5,j},\lambda_{6,j}$.
This argument uses the displayed coordinate
maps even when quotient characters repeat;
it requires no assumption about distinct
quotient characters within one factor.

\emph{Every comparison between different factors.}
We now prove that for $j\ne l$ there is no
source-intertwining isomorphism between such
subspaces of $Q_j$ and $Q_l$.
For a source operator acting by scalars $a$ and
$b$ on its two sides, an intertwiner $F$ satisfies
$(a-b)F=0$. A nonzero $F$ therefore forces $a=b$.
Apply this identity to the exact table (F.28.10).

A boundary-character subspace has source-$0$
scalar one. Every subspace contained in $W_l$
and every deleted-meridian-character subspace
has source-$0$ scalar $r_l\ne1$. This excludes
every comparison of a boundary character with
either of those types.

Two subspaces contained in $W_j$ and $W_l$
would force $r_j=r_l$ and $t_j=t_l$ through
sources $0$ and $7$. This contradicts (F.28.13).
For two boundary characters, source $7$ forces
$t_j=t_l$ and source $2$ then forces $r_j=r_l$,
giving the same contradiction. For the same
deleted type $4$ or $5$, sources $0$ and $7$
again determine $(r_j,t_j)$. For type $6$ the
source-$7$ value is $ct_j$, and the common
nonzero $c$ cancels, with the same result.

For two different deleted types, type $4$ versus
type $5$ is excluded by source $8$, whose scalars
are $c$ and $1$. Type $5$ versus type $6$ is
excluded by source $9$, with the same two values.
For type $4$ in factor $j$ versus type $6$ in
factor $l$, source $8$ would force
$t_l=c^{-1}$. Since $c\in k_1$, inverse Frobenius
then gives $t=c^{-1}$ and every $t_j=c^{-1}$.
Source $9$ would consequently compare $c$ with
$1$, again impossible because $c\ne1$.

Finally, a subspace of $W_j$ versus type $4$
in any factor is excluded by source $8$, acting
as $1$ and $c$. A subspace of $W_j$ versus
type $5$ is excluded by source $9$.
For type $6$ in factor $l$, source $8$ would
force $t_l=1$. Inverse Frobenius gives $t=1$
and all $t_j=1$. Source $7$ would then compare
$1$ on $W_j$ with $c$ on type $6$, a contradiction.
The inverse of an intertwining isomorphism
also intertwines, so these comparisons cover
both orientations of every pair of types.

\emph{The exact restriction kernel and affine image.}
Let
\[
 \Phi:\bigoplus_{j=0}^{m-1}Q_j
        \longrightarrow H^1_{\rm cont}(K,k_{0,\chi|K}),
        \qquad ([d_j])_j\longmapsto\sum_j[d_j|K].
                                                               \tag{F.28.14}
\]
This commutes with every actual source operator
$[d]\mapsto[d\circ B_l]$.
Each individual restriction $Q_j\to H^1(K,k_{0,\chi|K})$
is injective by F.24: its character on $R$ is
the same nontrivial $C$, and its coefficient
field is finite.

Suppose $\ker\Phi$ were nonzero, and choose a
nonzero invariant subspace $T$ inside it with
smallest dimension. Each nonzero coordinate
projection $p_j|T$ is injective and identifies
$T$ with an invariant subspace of $Q_j$ having
no proper nonzero invariant subspace. Two
nonzero projections would give the explicit
source-intertwining isomorphism
\[
                  p_l\circ(p_j|T)^{-1}:p_j(T)\longrightarrow p_l(T).
                                                               \tag{F.28.15}
\]
Its inverse is $p_j\circ(p_l|T)^{-1}$.
The comparisons just proved exclude this for
$j\ne l$. Thus $T$ has only one nonzero coordinate
projection and is contained in one $Q_j$.
On that factor $\Phi$ is the injective individual
restriction, a final contradiction. This proves
the summed $H^1$ restriction injective.

Consequently zero summed restriction of cocycles
forces every cocycle to be principal. As all
characters agree on $K$ and $C\ne1$, the exact
kernel is
\[
 \left\{\bigl(b_j(\chi_j-1)\bigr)_j:
                 b_j\in k_0,\ \sum_jb_j=0\right\}.
                                                               \tag{F.28.16}
\]
The converse inclusion follows by direct restriction.
Its dimension is $m-1$ on the $7m$ continuous
cocycle coordinates, proving rank $6m+1$.

A cocycle in the trace kernel of F.26.2 gives
the tuple $(d^{q_1^j})_j$ with zero summed
restriction. It is therefore principal, and
its parameter has trace zero because $C\ne1$.
Conversely every such principal cocycle has
zero trace restriction. The precise trace
kernel is thus $\mathscr P_\chi$ of F.26.3.
Every proper affine marginal image would give
a nonprincipal trace-kernel element by the
actual conjugation and coefficient homomorphism
of F.26.1--12. Those maps retain the original
scalar character, so they preserve (F.28.1)
and the fields (F.28.3). The proved kernel
equality excludes such an element. This
proves (F.28.2) for the remaining negative
case, completing all cases considered above.

\emph{Integral evidence and sign transport.}
The symbolic producer applies exactly (F.28.4)
to the retained F.24 universal data and records
all eleven full matrices, all eight seed rows,
all basis maps and all extension blocks.
A separately structured verifier imports
neither the producer nor an algebra package.
It evaluates all nineteen literal F.23 source
and seed braid words by generic affine
multiplication over the integral Laurent ring
$A$. It checks both basis inverses, the
principal-quotient identity after clearing
only $\delta_7$, the invariant inclusion,
each diagonal quotient coordinate, every
entry of (F.28.10), its exact source labels
and the actual seed evaluations (F.28.11).
It records $781$ checks. All coefficient
identities have integer coefficients; the
only nonmonomial localization is at $1-rst$,
whose nonvanishing was already proved.
The evidence manifest binds these programs,
outputs and inputs to the exact statement.

For a positive representation, precomposition
with the actual continuous isomorphism
$\mathcal B_+:V_-\to V_+$ sends the displayed
$Y_i^+$ character to precisely the negative
character just treated. This map and its
inverse carry $\mathcal K_-$ and $R_-$ onto
their positive counterparts by F.23.13.
Equality of the two negative images therefore
gives equality of the two original positive
images through the specified maps, proving
(F.28.2) for both signs.
\end{proof}
